\bmtchapitre{Formulaire}{Retrouver une formule avec ses conditions d’utilisation.}{Aide à la révision}
\section{Analyse}
Tangente : $y=f'(a)(x-a)+f(a)$, si $f$ est dérivable en $a$.
Pour $u,v$ dérivables : $(uv)'=u'v+uv'$, $(u/v)'=(u'v-uv')/v^2$ si $v\ne0$.
$\sqrt{x}$ a pour dérivée $1/(2\sqrt{x})$ sur $]0;+\infty[$.
Asymptote $y=ax+b$ : $f(x)-(ax+b)\to0$ en l’infini considéré.
\section{Algèbre}
Second degré ($a\ne0$) : $\Delta=b^2-4ac$, racines $(-b\pm\sqrt\Delta)/(2a)$ si $\Delta\ge0$.
Arithmétique : $u_n=u_p+(n-p)r$; somme de $k$ termes $k(\text{premier}+\text{dernier})/2$.
Géométrique ($q\ne0$) : $u_n=u_pq^{n-p}$; somme $u_p(1-q^k)/(1-q)$ si $q\ne1$, $ku_p$ si $q=1$. Si $q=0$, traiter le terme initial séparément.
Pour deux entiers positifs : $\pgcd(a,b)\ppcm(a,b)=ab$.
\section{Géométrie}
Barycentre : $\sum\alpha_i\vect{MA_i}=s\vect{MG}$ avec $s=\sum\alpha_i\ne0$.
$\sum\alpha_i MA_i^2=sMG^2+\sum\alpha_i GA_i^2$.
Al-Kashi : $BC^2=AB^2+AC^2-2AB\,AC\cos\widehat A$.
Produit vectoriel en repère orthonormé direct :
\[
(u_2v_3-u_3v_2\,;\,u_3v_1-u_1v_3\,;\,u_1v_2-u_2v_1).
\]
$\cos(a+b)=\cos a\cos b-\sin a\sin b$;
$\sin(a+b)=\sin a\cos b+\cos a\sin b$.
$\sin2x=2\sin x\cos x$, $\cos2x=2\cos^2x-1$.
\section{Données et probabilités}
Taux : $(V_f-V_i)/V_i$, avec $V_i>0$. Coefficients successifs : produit des $1+t$.
Pour $N=\sum n_i>0$ : $\bar x=\sum n_ix_i/N$, $V=\sum n_i(x_i-\bar x)^2/N$, $\sigma=\sqrt V$.
$A_n^p=n!/(n-p)!$, $\binom np=n!/[p!(n-p)!]$, $0\le p\le n$.
Équiprobabilité finie : $\Prob(A)=|A|/|\Omega|$, $\Omega\ne\varnothing$.
$\Prob(\overline A)=1-\Prob(A)$ et $\Prob(A\cup B)=\Prob(A)+\Prob(B)-\Prob(A\cap B)$.
